% Chapter 6: Case Study — this pipeline
\chapter{Case Study: This Article's Production Pipeline}
\label{ch:casestudy}

\section{Overview}

This article was generated by the \texttt{uoh-sqak-ex03} CrewAI pipeline, demonstrating
the very patterns it describes. The pipeline runs as follows:

\begin{enumerate}
  \item \textbf{Research phase}: ResearcherAgent queries DuckDuckGo with structured search
        terms, synthesises results into \texttt{workspace/research\_notes.md}.
  \item \textbf{Writing phase}: WriterAgent reads the research notes and produces six
        Markdown chapters, including the Hebrew BiDi chapter.
  \item \textbf{Editing phase}: EditorAgent enforces citation consistency, verifies table
        syntax, and marks mathematical relationships with \texttt{[NEEDS\_FORMULA]}.
  \item \textbf{LaTeX phase}: LaTeXAgent converts edited Markdown to \texttt{.tex}, generates
        the comparison chart via \texttt{ChartGeneratorTool}, and runs four LuaLaTeX passes.
\end{enumerate}

\section{Metrics}

\begin{longtable}{@{}lll@{}}
\toprule
\textbf{Metric} & \textbf{Value} & \textbf{Notes} \\
\midrule
\endfirsthead
\toprule
\textbf{Metric} & \textbf{Value} & \textbf{Notes} \\
\midrule
\endhead
\bottomrule
\endfoot
Pipeline agents   & 4         & Sequential CrewAI Process \\
Tasks defined     & 4         & 1 per agent \\
LaTeX passes      & 4         & lualatex → biber → lualatex → lualatex \\
Python files      & $\leq$16  & Each $\leq$150 lines (rule R7) \\
Test coverage     & $\geq$85\%& Enforced via pytest-cov \\
\bottomrule
\caption{Pipeline metrics}
\label{tab:metrics}
\end{longtable}

\section{Lessons Learned}

The file-based Skill layer proved particularly valuable: injecting \texttt{SKILL.md}
content into each agent's backstory allows rapid iteration on agent behaviour without
code changes \parencite{CrewAI2024}. The \texttt{ApiGatekeeper} singleton prevented
duplicate rate-limit logic across tools, a common source of bugs in multi-agent systems.

The human-in-the-loop approval gate (\texttt{ArticleSDK.approve\_markdown()}) between
Editor and LaTeX Producer prevented three compilation errors during development by allowing
manual correction of malformed Markdown tables before the expensive LaTeX pass.
